% !TEX root = ../../ctfp-print.tex

\lettrine[lhang=0.17]{モ}{ナドについてプログラマーに話すと}、きっとエフェクトの話になるだろう。数学者にとって、モナドは\newterm{代数}（algebra）に関するものだ。代数については後で話す---プログラミングにおいて重要な役割を果たすのだが---まず、モナドとの関係について少し直観を与えたい。今のところ、やや大雑把な議論だが、辛抱してほしい。

代数とは、式を作り、操作し、評価することについてだ。式は演算子を使って構築される。この単純な式を考えてみよう：
\[x^2 + 2 x + 1\]
この式は $x$ のような変数と $1$ や $2$ のような定数を、プラスや掛け算のような演算子で結びつけて形成されている。プログラマーとして、式を木として考えることが多い。

\begin{figure}[H]
  \centering
  \includegraphics[width=0.3\textwidth]{images/exptree.png}
\end{figure}

\noindent
木はコンテナなので、より一般的に言えば、式は変数を格納するコンテナだ。圏論では、コンテナを自己関手として表現する。変数 $x$ に型 $a$ を割り当てると、式は型 $m\ a$ を持ち、ここで $m$ は式の木を構築する自己関手だ。（自明でない分岐式は通常、再帰的に定義された自己関手を使って作られる。）

式に対して実行できる最も一般的な操作は何だろうか？それは代入だ：変数を式に置き換えること。例えば、先の例では $X$ を $y - 1$ に置き換えて次を得られる：
\[(y - 1)^2 + 2 (y - 1) + 1\]
何が起きたかというと：型 $m\ a$ の式を取り、型 $a \to m\ b$ の変換を適用した（$b$ は $y$ の型を表す）。結果は型 $m\ b$ の式だ。明示すると：
\[m\ a \to (a \to m\ b) \to m\ b\]
そう、これはモナドのバインドのシグネチャだ。

これが動機付けだった。では、モナドの数学に入ろう。数学者はプログラマーとは異なる記法を使う。自己関手には文字 $T$ を、ギリシャ文字として \code{join} には $\mu$ を、\code{return} には $\eta$ を好んで使う。\code{join} と \code{return} はどちらも多相関数なので、自然変換に対応すると推測できる。

したがって、圏論ではモナドは、$\mu$ と $\eta$ という一対の自然変換を備えた自己関手 $T$ として定義される。

$\mu$ は関手 $T^2$ の二乗から $T$ への自然変換だ。この二乗は単に関手を自身と合成したもの、$T \circ T$ だ（この種の二乗は自己関手に対してのみ行える）。
\[\mu \Colon T^2 \to T\]
この自然変換の対象 $a$ における成分は射：
\[\mu_a \Colon T (T a) \to T a\]
であり、$\Hask$ では、これは直接 \code{join} の定義に翻訳される。

$\eta$ は恒等関手 $I$ と $T$ の間の自然変換だ：
\[\eta \Colon I \to T\]
$I$ の対象 $a$ への作用は単に $a$ なので、$\eta$ の成分は射によって与えられる：
\[\eta_a \Colon a \to T a\]
これは直接 \code{return} の定義に翻訳される。

これらの自然変換はいくつかの追加の法則を満たさなければならない。その見方の一つは、これらの法則によって自己関手 $T$ のクライスリ圏を定義できるというものだ。$a$ と $b$ の間のクライスリ矢は $a \to T b$ という射として定義されることを思い出してほしい。二つのそのような矢の合成（添字 $T$ を付けた丸で書く）は $\mu$ を使って実装できる：
\[g \circ_T f = \mu_c \circ (T g) \circ f\]
ここで
\begin{gather*}
  f \Colon a \to T b \\
  g \Colon b \to T c
\end{gather*}
ここで $T$ は関手なので、射 $g$ に適用できる。Haskell 記法でこの式を認識するのは簡単かもしれない：

\src{snippet01}
または、成分では：

\src{snippet02}
代数的解釈の観点では、二つの連続した代入を合成しているだけだ。

クライスリ矢が圏を形成するためには、合成が結合的であり、$\eta_a$ が $a$ におけるクライスリの恒等矢であることが必要だ。この要求は $\mu$ と $\eta$ のモナド則に翻訳できる。しかし、これらの法則を導出する別の方法があり、それによってモノイドの法則のように見えるようになる。実際、$\mu$ はしばしば\emph{乗算}と呼ばれ、$\eta$ は\emph{単位}と呼ばれる。

大まかに言えば、結合律は、$T$ の三乗 $T^3$ を $T$ に還元する二つの方法が同じ結果を与えなければならないと述べている。二つの単位法則（左と右）は、$\eta$ を $T$ に適用し、$\mu$ で還元すると $T$ が戻ってくるという。

自然変換と関手を合成しているので少し複雑だ。水平合成について少し復習が必要だ。例えば、$T^3$ は $T$ が $T^2$ の後に合成されたものとして見られる。これに二つの自然変換の水平合成を適用できる：
\[I_T \circ \mu\]

\begin{figure}[H]
  \centering
  \includegraphics[width=0.4\textwidth]{images/assoc1.png}
\end{figure}

\noindent
そして $T \circ T$ を得る。これをさらに $\mu$ を適用して $T$ に還元できる。$I_T$ は $T$ から $T$ への恒等自然変換だ。この種の水平合成 $I_T \circ \mu$ の記法は $T \circ \mu$ と略されることが多い。この記法は明確だ。なぜなら、関手を自然変換と合成することは意味がないので、この文脈では $T$ は $I_T$ を意味するに違いないからだ。

\noindent
（自己）関手圏 ${[}\cat{C}, \cat{C}{]}$ の図式を描くこともできる：

\begin{figure}[H]
  \centering
  \includegraphics[width=0.3\textwidth]{images/assoc2.png}
\end{figure}

\noindent
あるいは、$T^3$ を $T^2 \circ T$ の合成として扱い、$\mu \circ T$ を適用することもできる。結果もまた $T \circ T$ であり、これも $\mu$ を使って $T$ に還元できる。二つのパスが同じ結果を生むことを要求する。

\begin{figure}[H]
  \centering
  \includegraphics[width=0.3\textwidth]{images/assoc.png}
\end{figure}

\noindent
同様に、水平合成 $\eta \circ T$ を恒等関手 $I$ が $T$ の後にある合成に適用して $T^2$ を得ることができる。これは $\mu$ を使って還元できる。結果は、恒等自然変換を直接 \code{T} に適用した場合と同じであるべきだ。そして、類推によって $T \circ \eta$ についても同じことが成り立つべきだ。

\begin{figure}[H]
  \centering
  \includegraphics[width=0.4\textwidth]{images/unitlawcomp-1.png}
\end{figure}

\noindent
これらの法則がクライスリ矢の合成が確かに圏の法則を満たすことを保証することを自分で確認できるだろう。

モナドとモノイドの類似性は際立っている。乗算 $\mu$、単位 $\eta$、結合律、単位律がある。しかし、モノイドの以前の定義はモナドをモノイドとして説明するには狭すぎる。そこで、モノイドの概念を一般化しよう。

\section{モノイダル圏}

モノイドの慣例的な定義に戻ろう。それは二項演算と単位と呼ばれる特別な元を持つ集合だ。Haskell では、これを型クラスとして表現できる：

\src{snippet03}
二項演算 \code{mappend} は結合的かつ単位的（つまり、単位 \code{mempty} による乗算は何もしない）でなければならない。

Haskell では、\code{mappend} の定義はカリー化されていることに注意しよう。これは \code{m} の各元を関数にマッピングすると解釈できる：

\src{snippet04}
この解釈が、モノイドを単一対象圏として定義する---ここで自己射 \code{(m -> m)} がモノイドの元を表す---ことへつながる。しかし、カリー化が Haskell に組み込まれているので、乗算の別の定義から始めることもできる：

\src{snippet05}
ここでは、デカルト積 \code{(m, m)} が乗算される対の源になる。

この定義は別の一般化への道を示唆している：デカルト積を圏論的積に置き換えること。積が大域的に定義されている圏から始め、そこで対象 \code{m} を選び、乗算を射として定義できる：
\[\mu \Colon m\times{}m \to m\]
ただし一つ問題がある：任意の圏では対象の内部を覗けないので、単位元をどのように選ぶのか？ここにトリックがある。元の選択はシングルトン集合からの関数と等価であることを思い出してほしい。Haskell では、\code{mempty} の定義を関数に置き換えることができる：

\src{snippet06}
シングルトンは $\Set$ における終対象なので、この定義を終対象 $t$ を持つ任意の圏に一般化するのが自然だ：
\[\eta \Colon t \to m\]
これにより元について話すことなく「単位」元を選ぶことができる。

単一対象圏としてのモノイドの以前の定義とは異なり、ここではモノイド則は自動的に満たされない---明示的に課す必要がある。しかし、それらを定式化するためには、基礎となる圏論的積自体のモノイダル構造を確立しなければならない。まず Haskell でモノイダル構造がどのように機能するかを振り返ろう。

結合律から始めよう。Haskell では、対応する等式的法則は：

\src{snippet07}
これを他の圏に一般化する前に、個々の変数への作用から抽象化して、関数（射）の等式として書き直す必要がある---言い換えれば、ポイントフリー記法を使う必要がある。デカルト積が双関手であることを知っているので、左辺を次のように書ける：

\src{snippet08}
右辺を：

\src{snippet09}
これはほぼ望むものだ。残念ながら、デカルト積は厳密には結合的でない---\code{(x, (y, z))} は \code{((x, y), z)} とは異なる---なのでポイントフリーで単純に書くことはできない：

\src{snippet10}
一方、対の二つのネストは同型だ。それらの間を変換する\newterm{結合子}（associator）と呼ばれる可逆関数がある：

\src{snippet11}
結合子の助けを借りて、\code{mu} のポイントフリー結合律を書ける：

\src{snippet12}
単位律にも同様のトリックを適用できる。新しい記法では次の形を取る：

\src{snippet13}
これは次のように書き直せる：

\src{snippet14}
同型 \code{lambda} と \code{rho} はそれぞれ\newterm{左単位子}（left unitor）と\newterm{右単位子}（right unitor）と呼ばれる。これらは単位 \code{()} が同型を除いてデカルト積の恒等元であるという事実を証明する：

\src{snippet15}

\src{snippet16}
したがって、単位律のポイントフリー版は：

\src{snippet17}
デカルト積自体が型の圏においてモノイダル乗算のように振る舞うという事実を使って、\code{mu} と \code{eta} のポイントフリーのモノイド則を定式化した。ただし、デカルト積の結合律と単位律は同型を除いてのみ有効であることに注意してほしい。

これらの法則は、積と終対象を持つ任意の圏に一般化できることがわかる。圏論的積は確かに同型を除いて結合的であり、終対象は単位であり、これも同型を除いてだ。結合子と二つの単位子は自然同型だ。法則は可換図式で表現できる。

\begin{figure}[H]
  \centering
  \includegraphics[width=0.5\textwidth]{images/assocmon.png}
\end{figure}

\noindent
積は双関手なので、射の対を持ち上げることができる---Haskell ではこれは \code{bimap} を使って行われた。

ここで止まって、圏論的積と終対象を持つ任意の圏の上にモノイドを定義できると言えるかもしれない。対象 $m$ と、モノイド則を満たす二つの射 $\mu$ と $\eta$ を選ぶことができる限り、モノイドを持つ。しかし、それよりも良いことができる。$\mu$ と $\eta$ の法則を定式化するために本格的な圏論的積は必要ない。積は射影を使う普遍的構成によって定義されることを思い出してほしい。モノイド則の定式化では射影を一切使っていない。

積のような振る舞いをするが積ではない双関手を\newterm{テンソル積}（tensor product）と呼び、しばしば中置演算子 $\otimes$ で表記される。一般的なテンソル積の定義は少し扱いにくいが、心配する必要はない。その性質を列挙するだけにしよう---最も重要なのは同型を除いた結合性だ。

同様に、対象 $t$ が終対象である必要もない。その終対象としての性質---すなわち、任意の対象からそれへの一意な射の存在---を一切使っていない。必要なのは、テンソル積とうまく協調して機能することだ。つまり、やはり同型を除いてテンソル積の単位であることを求める。まとめよう：

\newterm{モノイダル圏}（monoidal category）は、テンソル積と呼ばれる双関手を備えた圏 $\cat{C}$：
\[\otimes \Colon \cat{C}\times{}\cat{C} \to \cat{C}\]
と、単位対象と呼ばれる特別な対象 $i$、および結合子、左単位子、右単位子とそれぞれ呼ばれる三つの自然同型から成る：
\begin{align*}
  \alpha_{a b c} & \Colon (a \otimes b) \otimes c \to a \otimes (b \otimes c) \\
  \lambda_a      & \Colon i \otimes a \to a                                   \\
  \rho_a         & \Colon a \otimes i \to a
\end{align*}
（四重テンソル積を単純化するための整合条件もある。）

重要なのは、テンソル積が多くの身近な双関手を記述することだ。特に、積、余積、そしてすぐに見るように自己関手の合成（および Day 畳み込みのようないくつかのより難解な積）に対して機能する。モノイダル圏は豊かな圏の定式化において本質的な役割を果たす。

\section{モノイダル圏におけるモノイド}

これでモノイダル圏のより一般的な設定でモノイドを定義する準備ができた。対象 $m$ を選ぶことから始める。テンソル積を使って $m$ の冪を形成できる。$m$ の二乗は $m \otimes m$ だ。$m$ の三乗を形成する二つの方法があるが、それらは結合子を通じて同型だ。$m$ のより高い冪についても同様だ（そこで整合条件が必要になる）。モノイドを形成するには、二つの射を選ぶ必要がある：
\begin{align*}
  \mu  & \Colon m \otimes m \to m \\
  \eta & \Colon i \to m
\end{align*}
ここで $i$ はテンソル積の単位対象だ。

\begin{figure}[H]
  \centering
  \includegraphics[width=0.4\textwidth]{images/monoid-1.jpg}
\end{figure}

\noindent
これらの射は結合律と単位律を満たさなければならず、次の可換図式で表現できる：

\begin{figure}[H]
  \centering
  \includegraphics[width=0.5\textwidth]{images/assoctensor.jpg}
\end{figure}

\begin{figure}[H]
  \centering
  \includegraphics[width=0.5\textwidth]{images/unitmon.jpg}
\end{figure}

\noindent
テンソル積は双関手でなければならないことが本質的だ。なぜなら、$\mu \otimes \id$ や $\eta \otimes \id$ のような積を形成するために射の対を持ち上げる必要があるからだ。これらの図式は、圏論的積に対する以前の結果の直接的な一般化に過ぎない。

\section{モノイドとしてのモナド}

モノイダル構造は予期しない場所に現れる。そのような場所の一つが関手圏だ。少し目を細めると、関手の合成を乗算の一形式として見ることができるかもしれない。問題は、任意の二つの関手を合成できるわけではない---一方の目標圏が他方の源圏でなければならない。これは射の合成の通常の規則に過ぎない---そして知っての通り、関手は圏 $\Cat$ における射だ。しかし、自己射（同じ対象にループバックする射）が常に合成可能であるのと同様に、自己関手も常に合成可能だ。任意の圏 $\cat{C}$ について、$\cat{C}$ から $\cat{C}$ への自己関手は関手圏 ${[}\cat{C}, \cat{C}{]}$ を形成する。その対象は自己関手であり、射はそれらの間の自然変換だ。この圏から任意の二つの対象、例えば自己関手 $F$ と $G$ を取り、第三の対象 $F \circ G$---それらの合成である自己関手---を生成できる。

自己関手の合成はテンソル積の良い候補だろうか？まず、それが双関手であることを確立しなければならない。射の対---ここでは自然変換---を持ち上げるために使えるか？テンソル積に対する \code{bimap} の類似物のシグネチャはこのようになる：
\[\mathit{bimap} \Colon (a \to b) \to (c \to d) \to (a \otimes c \to b \otimes d)\]
対象を自己関手に、矢を自然変換に、テンソル積を合成に置き換えると：
\[(F \to F') \to (G \to G') \to (F \circ G \to F' \circ G')\]
これは水平合成の特別なケースとして認識できるかもしれない。

\begin{figure}[H]
  \centering
  \includegraphics[width=0.4\textwidth]{images/horizcomp.png}
\end{figure}

\noindent
また、自己関手の合成---新しいテンソル積---の恒等元として機能できる恒等自己関手 $I$ も手元にある。さらに、関手の合成は結合的だ。実際、結合律と単位律は厳格であり---結合子や二つの単位子は必要ない。したがって、自己関手は関手合成をテンソル積として\newterm{厳格なモノイダル圏}（strict monoidal category）を形成する。

この圏におけるモノイドは何だろうか？それは対象---すなわち自己関手 $T$---と二つの射---すなわち自然変換だ：
\begin{gather*}
  \mu \Colon T \circ T \to T \\
  \eta \Colon I \to T
\end{gather*}
それだけでなく、ここにモノイド則がある：

\begin{figure}[H]
  \centering
  \includegraphics[width=0.4\textwidth]{images/assoc.png}
\end{figure}

\begin{figure}[H]
  \centering
  \includegraphics[width=0.5\textwidth]{images/unitlawcomp.png}
\end{figure}

\noindent
これらはまさに以前見たモナド則だ。これで Saunders Mac Lane の有名な引用が理解できる：

\begin{quote}
  つまるところ、モナドは自己関手圏におけるモノイドに過ぎない。
\end{quote}
関数型プログラミングのカンファレンスでTシャツに書かれているのを見たことがあるかもしれない。

\section{随伴からのモナド}

\hyperref[adjunctions]{随伴}\footnote{随伴については第18章を参照。}$L \dashv R$ は二つの圏 $\cat{C}$ と $\cat{D}$ の間を行き来する一対の関手だ。それらを合成する二つの方法があり、二つの自己関手 $R \circ L$ と $L \circ R$ が生まれる。随伴に従って、これらの自己関手は単位と余単位と呼ばれる二つの自然変換を通じて恒等関手と関連している：
\begin{gather*}
  \eta \Colon I_{\cat{D}} \to R \circ L \\
  \varepsilon \Colon L \circ R \to I_{\cat{C}}
\end{gather*}
すぐに気づくのは、随伴の単位はモナドの単位とまったく同じように見えることだ。実は、自己関手 $R \circ L$ は確かにモナドだ。$\eta$ と組み合わせる適切な $\mu$ を定義するだけでよい。それは自己関手の二乗と自己関手自体の間の自然変換であり、随伴関手の観点では：
\[R \circ L \circ R \circ L \to R \circ L\]
そして確かに、中間の $L \circ R$ を崩すために余単位を使うことができる。$\mu$ の正確な式は水平合成によって与えられる：
\[\mu = R \circ \varepsilon \circ L\]
モナド則は随伴の単位と余単位が満たす恒等式と交換法則から従う。

Haskell では随伴から導かれるモナドをあまり見かけない。なぜなら、随伴は通常二つの圏を含むからだ。しかし、指数または関数対象の定義は例外だ。この随伴を形成する二つの自己関手はこれだ：
\begin{gather*}
  L z = z\times{}s \\
  R b = s \Rightarrow b
\end{gather*}
それらの合成を見慣れた状態モナドとして認識できるかもしれない：
\[R (L z) = s \Rightarrow (z\times{}s)\]
このモナドは Haskell で以前に見た：

\src{snippet18}
随伴も Haskell に翻訳しよう。左関手は積関手だ：

\src{snippet19}
右関手は読み手関手だ：

\src{snippet20}
それらは随伴を形成する：

\src{snippet21}
積関手の後に読み手関手を合成することが状態関手と等価であることを簡単に確認できる：

\src{snippet22}
予想通り、随伴の \code{unit} は状態モナドの \code{return} 関数と等価だ。\code{counit} はその引数に作用する関数を評価することで機能する。これは関数 \code{runState} のカリー化されていないバージョンとして認識できる：

\src{snippet23}
（\code{counit} では対に作用するため、カリー化されていない。）

これで自然変換 $\mu$ の成分として状態モナドの \code{join} を定義できる。そのために、三つの自然変換の水平合成が必要だ：
\[\mu = R \circ \varepsilon \circ L\]
言い換えれば、余単位 $\varepsilon$ を読み手関手の一層を越えて忍び込ませる必要がある。コンパイラが \code{Reader} 関手ではなく \code{State} 関手のものを選んでしまうので、直接 \code{fmap} を呼び出すことはできない。しかし、読み手関手の \code{fmap} は単に左からの関数合成であることを思い出してほしい。そこで関数合成を直接使う。

まず \code{State} データコンストラクタを剥がして、\code{State} 関手の内部の関数を露出させる必要がある。これは \code{runState} を使って行う：

\src{snippet24}
次にそれを余単位（\code{uncurry runState} で定義される）と左合成する。最後に \code{State} データコンストラクタで包み直す：

\src{snippet25}
これは確かに \code{State} モナドの \code{join} の実装だ。

随伴からモナドが生まれるだけでなく、逆もまた真であることがわかる：すべてのモナドは二つの随伴関手の合成に因数分解できる。ただし、そのような因数分解は一意ではない。

他の自己関手 $L \circ R$ については次のセクションで話す。
